\section{Heuristic motivation: Ricci flow parallels}
\label{app:RicciFlowPenrose}

This appendix discusses heuristic arguments inspired by Ricci flow that motivate the rigorous approach.

\subsection{Parallel with Hamilton's program}

Perelman's resolution of the Poincar\'e conjecture \cite{perelman2002,perelman2003} proceeds via geometric flows moving objects toward canonical forms, monotonicity formulas controlling the evolution, no-local-collapsing theorems, surgery past singularities, and convergence to special geometries.

For the spacetime Penrose inequality, analogous structures appear: the $\theta^+$-flow evolves trapped surfaces toward MOTS; a monotonicity formula controls mass and area; the dominant energy condition plays the role of positive curvature; and the final state corresponds to the horizon.
\begin{equation}
\frac{d}{dt}\mathrm{Area}(S_t) = -\int_{S_t} H \cdot \theta^+ \, dA.
\label{eq:area-evo-sign-app}
\end{equation}
This is \textbf{not generally monotone} (unlike Huisken's IMCF), but the sign structure reveals deep information about the trapped region geometry.

\subsection{The Perelman Entropy and Its Spacetime Analogue}

\subsubsection{Perelman's $\mathcal{W}$-Functional: Review}

In Riemannian geometry with Ricci flow $\partial_t g = -2\Ric$, Perelman introduced the entropy functional
\begin{equation}
\mathcal{W}(g, f, \tau) = \int_M \left[\tau\left(|\nabla f|^2 + R\right) + f - n\right]e^{-f} \, dV
\label{eq:perelman-W-app}
\end{equation}
where $f$ is a scalar field (the "entropy potential") and $\tau > 0$ is a scale parameter. The key properties are:
\begin{enumerate}[label=(\alph*)]
\item \textbf{Monotonicity:} $\frac{d}{d\tau}\mathcal{W} \geq 0$ under coupled flow $\partial_\tau f = -\Delta f + |\nabla f|^2 - R + \frac{n}{2\tau}$;
\item \textbf{Rigidity:} Equality holds if and only if $(M,g)$ is a gradient shrinking soliton;
\item \textbf{No-local-collapsing:} The $\mathcal{W}$-bound prevents degenerate blow-ups.
\end{enumerate}

\subsubsection{Spacetime Analogue: The Hawking-Geroch Entropy}

For a spacetime $(M^4, g_{\mu\nu})$ with spatial slice $(\Sigma^3, g, k)$, the natural entropy (in the time-symmetric case $k=0$) is the \textbf{Hawking quasi-local mass}:
\begin{equation}
m_H(S_t) 
= \sqrt{\frac{A(S_t)}{16\pi}}\left(1 - \frac{1}{16\pi}\int_{S_t} H^2 \, dA\right)
\label{eq:hawking-mass-app}
\end{equation}
where $S_t$ is a closed 2-surface in $\Sigma$.

\begin{proposition}[Hawking Mass and IMCF (Context)]
\label{prop:hawking-imcf-app}
In the time-symmetric case $k=0$, Geroch's computation shows that along inverse mean curvature flow $\partial_t S = H^{-1}\nu$ with $H>0$, the Hawking mass is non-decreasing under the hypothesis $R_g\ge 0$ (in the smooth setting).
\end{proposition}

\begin{proof}[Proof Sketch]
The evolution of area under $\dot{S} = H^{-1}\nu$ is $\frac{dA}{dt} = \int H^{-1} \cdot H \, dA = A(S_t)$, giving exponential growth. One then combines the Gauss equation for $S_t\subset (\Sigma,g)$, the evolution equation for $H$ along IMCF, and the hypothesis $R_g\ge 0$ to show $\frac{d}{dt}m_H(S_t)\ge 0$.

For the fully spacetime (non-time-symmetric) case $k\ne 0$, the correct monotone quantity and hypotheses involve additional terms (e.g. null expansions and suitable energy conditions), so we use this proposition only as motivation/analogy.
\end{proof}

\textbf{Problem for trapped surfaces:} IMCF requires $H > 0$, but trapped surfaces have $H + \tr_S k \leq 0$. In the "unfavorable" regime $\tr_S k < 0$, we can have $H > 0$ yet $\theta^+ \leq 0$---a purely spacelike positive-mean-curvature surface that is nonetheless trapped.

\subsection{A \texorpdfstring{$\theta^+$}{theta+}-Adjusted Entropy Functional (Heuristic Program)}

\begin{definition}[Spacetime Perelman Functional (proposal)]
\label{def:spacetime-perelman-app}
For a trapped surface $S \subset \Sigma$ evolving under the $\theta^+$-flow, define
\begin{equation}
\mathcal{P}(S, \phi) = \int_S \left[(\theta^+)^2 + |\nabla\phi|^2 + \phi \cdot R_S\right]e^{-\phi} \, dA
\label{eq:penrose-entropy-app}
\end{equation}
where $\phi: S \to \mathbb{R}$ is an auxiliary "entropy field" and $R_S$ is the Gauss curvature of $S$.
\end{definition}

\begin{theorem}[Monotonicity of $\mathcal{P}$ under Coupled Flow (formal computation)]
\label{thm:spacetime-monotonicity-app}
Let $S_t$ evolve under $\frac{\partial S}{\partial t} = -\theta^+(S)\nu$ (the $\theta^+$-flow), and let $\phi$ evolve under
\begin{equation}
\frac{\partial \phi}{\partial t} = -\Delta_S \phi + |\nabla \phi|^2 - (\theta^+)^2 + R_S.
\label{eq:phi-evolution-app}
\end{equation}
Assume $S_t$ remains smooth and closed for the time interval under consideration, and that all geometric quantities needed below are smooth.
If, in addition, one has the \emph{auxiliary coercivity/curvature hypotheses} required to control the lower-order terms in the $\theta^+$ evolution (see Remark~\ref{rem:P-functional-status-app}), then the following formal computation suggests that
\begin{equation}
\frac{d}{dt}\mathcal{P}(S_t, \phi_t) \geq 0
\label{eq:P-monotone-app}
\end{equation}
with equality only in the (formal) ``stationary'' situation where $\theta^+ \equiv 0$ and $\phi$ is constant (up to the usual gauge normalizations for weighted energies).
\end{theorem}

\begin{remark}[Status of the $\mathcal{P}$-functional computation]
\label{rem:P-functional-status-app}
The material in this subsection is included as a \emph{Ricci-flow-inspired heuristic} rather than as an input to the main proofs in this paper.
At present, we do \emph{not} supply a complete set of hypotheses under which \eqref{eq:P-monotone-app} can be established for the coupled system
\eqref{eq:phi-evolution-app} together with a well-posed geometric flow driven by $-\theta^+\nu$.
In particular, turning the computation into a theorem would require at least:
\begin{itemize}
\item a precise evolution equation for $\theta^+$ under the chosen flow and an identification of all lower-order terms (including those involving $k$ and ambient spacetime curvature), which are currently suppressed in the schematic Equation~\eqref{eq:theta-evolution-schematic-app};
\item a coercive estimate that controls these lower-order contributions and justifies the completion-of-squares step globally in time;
\item a proof of short-time existence and regularity for the coupled geometric/PDE system, or an appropriate weak formulation.
\end{itemize}
For the rigorous arguments establishing the Penrose inequality in this work, see the Jang--conformal--AMO pipeline developed in Sections~\ref{sec:Analysis} and~\ref{sec:Synthesis}.
\end{remark}

\begin{proof}
We present a \emph{formal} computation (i.e. ignoring several lower-order terms and regularity issues). Decompose the time derivative:
\begin{equation}
\frac{d}{dt}\mathcal{P} = \underbrace{\int_S \partial_t\left[(\theta^+)^2 + |\nabla\phi|^2 + \phi R_S\right]e^{-\phi} \, dA}_{I_1} + \underbrace{\int_S \left[(\theta^+)^2 + |\nabla\phi|^2 + \phi R_S\right]\partial_t(e^{-\phi} \, dA)}_{I_2}.
\end{equation}

\textbf{Step 1: Evolution of the weighted area form.}
Under normal velocity $V = -\theta^+$ (with respect to the outward unit normal $\nu$), the first variation formula gives
\begin{equation}
\partial_t(dA) = H\,V\, dA = -H\,\theta^+\, dA.
\end{equation}
Combined with $\partial_t(e^{-\phi}) = -(\partial_t\phi)e^{-\phi}$:
\begin{equation}
\partial_t(e^{-\phi} \, dA) = e^{-\phi}\left(-\partial_t\phi - H\theta^+\right)dA.
\end{equation}
Substituting the evolution \eqref{eq:phi-evolution-app}:
\begin{equation}
-\partial_t\phi - H\theta^+ = \Delta_S\phi - |\nabla\phi|^2 + (\theta^+)^2 - R_S - H\theta^+.
\end{equation}

\textbf{Step 2: Evolution of $\theta^+$ under the $\theta^+$-flow.}
We use only that, for smooth deformations, the principal part of the linearization of $\theta^+$ in the normal direction is elliptic on $S$ (a Laplace--Beltrami term). More precisely, one expects a schematic evolution of the form
\begin{equation}
\partial_t\theta^+ = -\Delta_S\theta^+ + \text{(lower order terms depending on } g,k,\text{ and ambient curvature)}.
\label{eq:theta-evolution-schematic-app}
\end{equation}
We will not attempt to identify all lower-order terms here.

\textbf{Step 3: Where DEC would enter.}
In a fully spacetime formulation, the matter term $T(\ell,\ell)$ (or equivalently $G(\ell,\ell)$ via Einstein's equations) appears in the null Raychaudhuri equation and in stability/variation formulas for null expansions. Under appropriate energy conditions one can obtain favorable-sign contributions. We do not use any specific inequality beyond this qualitative comment in the remainder of this formal computation.

\textbf{Step 4: Evolution of $(\theta^+)^2$.}
From \eqref{eq:theta-evolution-schematic-app} and keeping only the principal part explicitly:
\begin{align}
\partial_t(\theta^+)^2 &= 2\theta^+\partial_t\theta^+ = -2\theta^+\Delta_S\theta^+.
\end{align}
Integrating by parts on the closed surface $S$:
\begin{equation}
\int_S \partial_t(\theta^+)^2 \, e^{-\phi}dA = \int_S 2|\nabla\theta^+|^2 e^{-\phi}dA - \int_S 2\theta^+\nabla\theta^+\cdot\nabla\phi \, e^{-\phi}dA.
\label{eq:theta-sq-evol-app}
\end{equation}

\textbf{Step 5: Evolution of $|\nabla\phi|^2$.}
On an evolving surface, commuting $\partial_t$ with intrinsic covariant derivatives produces additional curvature and second fundamental form terms. We suppress these and focus only on the Bochner identity contribution that appears in the usual Perelman-style completion of squares.
\begin{align}
\partial_t|\nabla\phi|^2 &= 2\langle\nabla\phi, \nabla(\partial_t\phi)\rangle + 2\langle\nabla\phi, (\nabla\theta^+)\cdot\nabla\phi\rangle \\
&\approx 2\langle\nabla\phi, \nabla(-\Delta_S\phi + |\nabla\phi|^2 - (\theta^+)^2 + R_S)\rangle.
\end{align}
Expanding the first term:
\begin{align}
2\langle\nabla\phi, \nabla(-\Delta_S\phi)\rangle &= -2\langle\nabla\phi, \nabla\Delta_S\phi\rangle.
\end{align}
The Bochner identity on the 2-surface $S$ states:
\begin{equation}
\frac{1}{2}\Delta_S|\nabla\phi|^2 = |\nabla^2\phi|^2 + \langle\nabla\phi, \nabla\Delta_S\phi\rangle + \frac{R_S}{2}|\nabla\phi|^2.
\end{equation}
Therefore:
\begin{equation}
-2\langle\nabla\phi, \nabla\Delta_S\phi\rangle = -\Delta_S|\nabla\phi|^2 + 2|\nabla^2\phi|^2 + R_S|\nabla\phi|^2.
\end{equation}

\textbf{Step 6: Combine all terms.}
After integration by parts and collecting terms with the weight $e^{-\phi}$:
\begin{align}
\frac{d}{dt}\mathcal{P} &= \int_S e^{-\phi}\Bigg[2|\nabla\theta^+|^2 + 2|\nabla^2\phi|^2 + R_S|\nabla\phi|^2 \notag\\
&\quad - 2\theta^+\nabla\theta^+\cdot\nabla\phi + 4\langle\nabla\phi, \nabla|\nabla\phi|^2\rangle - 2\langle\nabla\phi, \nabla(\theta^+)^2\rangle \notag\\
&\quad + \left((\theta^+)^2 + |\nabla\phi|^2 + \phi R_S\right)\left(\Delta_S\phi - |\nabla\phi|^2 + (\theta^+)^2 - R_S - H\theta^+\right)\Bigg]dA.
\label{eq:P-evolution-full-app}
\end{align}

\textbf{Step 7: Complete the square (formal).}
The cross terms $-2\theta^+\nabla\theta^+\cdot\nabla\phi$ and $-2\langle\nabla\phi, \nabla(\theta^+)^2\rangle = -4\theta^+\nabla\theta^+\cdot\nabla\phi$ combine to give:
\begin{equation}
-6\theta^+\nabla\theta^+\cdot\nabla\phi = -6\theta^+\langle\nabla\theta^+, \nabla\phi\rangle.
\end{equation}
By Cauchy-Schwarz with parameter $\epsilon > 0$:
\begin{equation}
\left|6\theta^+\langle\nabla\theta^+, \nabla\phi\rangle\right| \leq 3\epsilon|\nabla\theta^+|^2 + \frac{3}{\epsilon}(\theta^+)^2|\nabla\phi|^2.
\end{equation}
Choosing $\epsilon = 2/3$:
\begin{equation}
-6\theta^+\langle\nabla\theta^+, \nabla\phi\rangle \geq -2|\nabla\theta^+|^2 - \frac{9}{2}(\theta^+)^2|\nabla\phi|^2.
\end{equation}
Substituting back, the $|\nabla\theta^+|^2$ terms cancel, leaving:
\begin{align}
\frac{d}{dt}\mathcal{P} &\geq \int_S e^{-\phi}\Bigg[2|\nabla^2\phi|^2 + \left(R_S - \frac{9}{2}(\theta^+)^2\right)|\nabla\phi|^2 + (\theta^+)^4 \notag\\
&\quad - H\theta^+(\theta^+)^2 + \text{lower order}\Bigg]dA.
\end{align}

\textbf{Step 8: Sign analysis for trapped surfaces (incomplete).}
For a trapped surface, $\theta^+ \leq 0$, so $(\theta^+)^4 \geq 0$. The term $-H\theta^+(\theta^+)^2 = -H(\theta^+)^3$:
\begin{itemize}
\item If $H \geq 0$: $-H(\theta^+)^3 \geq 0$ since $(\theta^+)^3 \leq 0$.
\item If $H < 0$: additional geometric control is needed; we do not claim a general sign.
\end{itemize}

By the Gauss--Bonnet theorem, $\int_S R_S \, dA = 8\pi\chi(S)$, so for $S\cong S^2$ one has $\int_S R_S\,dA=8\pi$. This is only an averaged statement and does not by itself control $\int_S R_S|\nabla\phi|^2$.

\textbf{Conclusion:} Making \eqref{eq:P-monotone-app} rigorous would require: (i) an exact evolution equation for $\theta^+$ under the chosen flow; (ii) sharp control of all lower-order terms (including those involving $k$); and (iii) a coercive inequality to dominate mixed terms after integration by parts. We do not claim such a theorem here. Instead, the (very) schematic analysis suggests at best a Gr\"onwall-type inequality of the form
\begin{equation}
\mathcal{P}(S_t, \phi_t) \geq e^{-C't}\mathcal{P}(S_0, \phi_0).
\end{equation}

\textbf{Refined monotonicity:} Achieving a genuine inequality $\frac{d}{dt}\mathcal{P} \geq 0$ would require an 
\emph{a priori} mechanism to dominate the mixed terms and the (suppressed) lower-order contributions. Any pointwise condition such as $R_S \gtrsim (\theta^+)^2$ would have to be derived from the flow and ambient geometry; we do not assume or prove such a condition here.
\end{proof}

\subsection{Consequences for Geometric Control}

The entropy bounds provide geometric control, but we must be precise about what they do and do not imply.

\begin{corollary}[Curvature Bounds from Entropy]
\label{cor:curvature-bounds-app}
Let $S_0$ be a smooth trapped surface with $\theta^+(S_0) = -\epsilon_0 < 0$ and area $A_0 = \mathrm{Area}(S_0)$. Suppose the entropy satisfies the exponential bound
\begin{equation}
\mathcal{P}(S_t, \phi_t) \geq e^{-C't}\mathcal{P}(S_0, \phi_0)
\end{equation}
from Theorem~\ref{thm:spacetime-monotonicity-app}. Then along the $\theta^+$-flow, the $L^2$ norm of the second fundamental form satisfies
\begin{equation}
\int_{S_t} |A|^2 \, dA \leq C_1 e^{C't}\mathcal{P}(S_0, \phi_0) + C_2 \mathrm{Area}(S_t)
\end{equation}
for constants $C_1, C_2$ depending on $\|k\|_{L^\infty}$ and $\|R_g\|_{L^\infty}$.
\end{corollary}

\begin{proof}
The Gauss equation gives, for a 2-surface $S\subset (M^3,g)$,
\begin{equation}
R_S = R_g - 2\Ric(\nu,\nu) + H^2 - |A|^2,
\end{equation}
where $\nu$ is the unit normal and $A$ is the second fundamental form of $S$ in $(M,g)$.
For a 2-surface, Gauss--Bonnet implies
\begin{equation}
\int_S R_S \, dA = 4\pi \chi(S) = 8\pi
\end{equation}
for $S \cong S^2$. Rearranging:
\begin{equation}
\int_S |A|^2 \, dA = \int_S (R_g - 2\Ric(\nu,\nu) + H^2) \, dA - 8\pi.
\end{equation}

The dominant energy condition enters through the constraint equations. In general,
$R_g + (\tr k)^2 - |k|^2 = 16\pi\mu$ and $\Div(k-(\tr k)g)=8\pi J$.
In particular, DEC ($\mu\ge |J|$) does not imply a pointwise lower bound for $R_g$ alone unless one imposes additional gauge/size assumptions on $k$.
For the present heuristic estimate, we simply assume $R_g$ is bounded below on $S$ in the weak sense
\begin{equation}
\int_S R_g\,dA \ge -C(K)\,\Area(S)
\end{equation}
for some constant $C(K)$ depending on an \emph{a priori} $L^\infty$ bound on $k$ and the geometry in the region swept out by the flow.

The mean curvature $H = \theta^+ - \tr_S k$ satisfies:
\begin{equation}
\int_S H^2 \, dA \leq 2\int_S (\theta^+)^2 \, dA + 2\int_S (\tr_S k)^2 \, dA \leq 2\int_S (\theta^+)^2 \, dA + 2K^2 \mathrm{Area}(S).
\end{equation}

The entropy $\mathcal{P}$ contains $\int_S (\theta^+)^2 e^{-\phi} dA$. If $\phi$ is bounded (which requires separate analysis), we obtain:
\begin{equation}
\int_S (\theta^+)^2 \, dA \leq e^{\|\phi\|_{L^\infty}} \mathcal{P}(S_t, \phi_t).
\end{equation}

Combining these estimates yields the stated bound.
\end{proof}

\begin{remark}[What This Does NOT Prove]
\label{rem:no-collapse-gap-app}
The above corollary controls the \emph{integrated} curvature $\|A\|_{L^2}^2$, but this does \textbf{not} directly imply:
\begin{enumerate}[label=(\roman*)]
\item \textbf{Injectivity radius bounds:} The classical Klingenberg theorem requires pointwise curvature bounds and applies to complete Riemannian manifolds, not embedded surfaces. For surfaces, the relevant estimate is via Gauss--Bonnet and isoperimetric inequalities, which control area but not injectivity radius directly.
\item \textbf{$C^{2,\alpha}$ regularity:} $L^2$ curvature bounds do not prevent point concentration. Higher regularity requires Schauder estimates on the flow equation.
\item \textbf{Long-time existence:} Curvature blow-up ($|A| \to \infty$ at a point) can occur even with bounded $L^2$ norm.
\end{enumerate}
A complete no-local-collapsing theorem would require additional estimates---this is an \textbf{open problem} for the $\theta^+$-flow.
\end{remark}

\begin{corollary}[Weak Area Control]
\label{cor:weak-area-app}
Under the hypotheses of Corollary~\ref{cor:curvature-bounds-app}, if additionally the flow converges to a MOTS $\mathcal{M}$ as $t \to T^* < \infty$, then
\begin{equation}
\mathrm{Area}(\mathcal{M}) \leq e^{C'T^*}\left(\mathrm{Area}(S_0) + C''\mathcal{P}(S_0, \phi_0)\right).
\end{equation}
This gives an \textbf{upper} bound on $\mathrm{Area}(\mathcal{M})$, not the lower bound needed for Penrose.
\end{corollary}

\begin{proof}
The area evolution under normal velocity $V = -\theta^+$ is:
\begin{equation}
\frac{d}{dt}\mathrm{Area}(S_t) = -\int_{S_t} H\theta^+ \, dA.
\end{equation}
Using $|H\theta^+| \leq |H||\theta^+| \leq \frac{1}{2}(H^2 + (\theta^+)^2)$:
\begin{equation}
\left|\frac{d}{dt}\mathrm{Area}(S_t)\right| \leq \frac{1}{2}\int_{S_t} (H^2 + (\theta^+)^2) \, dA \leq C(\mathcal{P}(S_t) + \mathrm{Area}(S_t)).
\end{equation}
\end{proof}
